\section*{Capítulo \ref{ch:b2:ligand-field} --- Teoría del campo de ligandos y color}

\begin{solution}{exo:b2:ligand-field:1}
$d^3$: $t_{2g}^3$, EECC $-1.2\Delta_o$. $d^8$: $t_{2g}^6e_g^2$, EECC $-2.4 + 1.2 = -1.2\Delta_o$.
\end{solution}

\begin{solution}{exo:b2:ligand-field:2}
\ce{[Fe(H2O)6]^2+}, $d^6$ de \omterm{def:b2:ligand-field:spin}{espín alto} $t_{2g}^4e_g^2$: cuatro. \ce{[Fe(CN)6]^4-}, de \omterm{def:b2:ligand-field:spin}{espín bajo},
$t_{2g}^6$: ninguno (\omterm{def:b2:diatomic-mos:magnetism}{diamagnético}).
\end{solution}

\begin{solution}{exo:b2:ligand-field:3}
600 nm corresponde a luz naranja; el complejo se ve azul (azul verdoso).
\end{solution}

\begin{solution}{exo:b2:ligand-field:4}
$\ce{Cl-} < \ce{H2O} < \ce{NH3} < \ce{CN-}$.
\end{solution}

\begin{solution}{exo:b2:ligand-field:5}
Agua: $\Delta_o < P$, \omterm{def:b2:ligand-field:spin}{espín alto}, $t_{2g}^3e_g^2$, cinco electrones desapareados. Cianuro: $\Delta_o > P$,
\omterm{def:b2:ligand-field:spin}{espín bajo}, $t_{2g}^5$, un electrón desapareado.
\end{solution}

\begin{solution}{exo:b2:ligand-field:6}
En el \ce{[CoCl4]^2-} tetraédrico el desdoblamiento es de unos $4/9$ del octaédrico, y
el cloruro es un ligando más débil que el agua: la banda d--d se desplaza a longitudes de onda mayores (del naranja
al rojo), y el complejo se ve azul. Los complejos tetraédricos también absorben con más intensidad,
de ahí el color intenso.
\end{solution}

\begin{solution}{exo:b2:ligand-field:7}
\ce{Ni^2+} es $d^8$. Plano-cuadrado: los cuatro orbitales inferiores contienen los ocho electrones apareados,
y $\mathrm d_{x^2-y^2}$ queda vacío: \omterm{def:b2:diatomic-mos:magnetism}{diamagnético}. Tetraédrico: $e^4t_2^4$, dos electrones desapareados
en el conjunto $t_2$.
\end{solution}

\begin{solution}{exo:b2:ligand-field:8}
$\tilde\nu = 1/(500 \times 10^{-7}\,\text{cm}) = \qty{20000}{cm^{-1}}$; $\Delta_o =
0.1196/(500 \times 10^{-9}) = \qty{2.39e5}{J/mol}$, \qty{239}{kJ/mol}.
\end{solution}

\begin{solution}{exo:b2:ligand-field:9}
\ce{Zn^2+} es $d^{10}$ y \ce{Sc^3+}, $d^0$: no es posible ninguna \omterm{def:b2:ligand-field:d-d}{transición d--d} (no hay nivel d vacío
para \ce{Zn^2+} ni electrón d para \ce{Sc^3+}), de modo que no hay absorción en el visible.
\end{solution}

\begin{solution}{exo:b2:ligand-field:10}
Co(III) $d^6$ + seis pares no enlazantes del amoníaco: 18 electrones. Doce llenan los seis \omterm{def:b2:diatomic-mos:bonding}{orbitales enlazantes},
y seis llenan los $t_{2g}$ no enlazantes (el amoníaco da un campo lo bastante fuerte para el \omterm{def:b2:ligand-field:spin}{espín bajo} con el
cobalto(III)); $e_g^*$ está vacío. Todos los electrones están apareados: \omterm{def:b2:diatomic-mos:magnetism}{diamagnético}.
\end{solution}

\begin{solution}{exo:b2:ligand-field:11}
\ce{CO} es un buen dador $\sigma$ y, sobre todo, un aceptor $\pi$: sus orbitales $\pi^*$ vacíos
rebajan el nivel $t_{2g}$ por \omterm{def:b2:ligand-field:pi-ligands}{retrodonación}, ensanchando $\Delta_o$. \ce{F-} es un dador $\pi$, que
eleva $t_{2g}$ y estrecha $\Delta_o$. El modelo de cargas puntuales solo ve la carga y se equivoca en el
orden.
\end{solution}

\begin{solution}{exo:b2:ligand-field:12}
Sin campo, las entalpías de hidratación seguirían una línea regular (iones que se contraen a lo largo de
la serie). El agua da acuocomplejos de \omterm{def:b2:ligand-field:spin}{espín alto} cuya EECC es nula para $d^0$, $d^5$, $d^{10}$
y máxima para $d^3$ y $d^8$: la estabilización adicional añade dos jorobas, con máximos cerca de
\ce{V^2+} y de \ce{Ni^2+}, por encima de la recta que pasa por \ce{Ca^2+}, \ce{Mn^2+} y \ce{Zn^2+}.
\end{solution}

\begin{solution}{pb:b2:ligand-field:1}
\textbf{1.} $d^3$.
\textbf{2.} $t_{2g}^3$, un electrón en cada uno de los tres orbitales, con espines paralelos.
\textbf{3.} No: tres electrones caben en $t_{2g}$ sin aparearse.
\textbf{4.} $-1.2\Delta_o$.
\textbf{5.} Tres.
\textbf{6.} La EECC es grande y los orbitales $e_g^*$, que apuntan hacia los ligandos, están
vacíos: retirar o añadir un ligando cuesta buena parte de esa estabilización.
\textbf{7.} Seis iones óxido en los vértices de un octaedro ligeramente distorsionado.
\textbf{8.} $1/(556 \times 10^{-7}) = \qty{17990}{cm^{-1}}$, unos \qty{1.80e4}{cm^{-1}}.
\textbf{9.} $0.1196/(556 \times 10^{-9}) = \qty{2.15e5}{J/mol}$: \qty{215}{kJ/mol}.
\textbf{10.} Amarillo verdoso (556 nm) y violeta (405 nm).
\textbf{11.} Rojo, con un poco de azul: el rojo intenso del rubí.
\textbf{12.} Un ion $d^3$ tiene más de una configuración excitada con un electrón en $e_g$;
la repulsión entre electrones les da energías distintas, y de ahí varias bandas (se trata en el
volumen del tercer año).
\textbf{13.} $0.1196/(620 \times 10^{-9}) = \qty{1.93e5}{J/mol}$, \qty{193}{kJ/mol}
(\qty{16130}{cm^{-1}}).
\textbf{14.} El rubí: su desdoblamiento es mayor.
\textbf{15.} Hacia el rojo anaranjado: ahora se absorbe el rojo, y se transmite el verde (con algo de azul).
\textbf{16.} En el berilo los sitios del cromo son algo mayores y los iones óxido están más lejos,
y el entorno es distinto: un campo más débil.
\textbf{17.} No: un ion $d^3$ conserva tres electrones desapareados en cualquier campo octaédrico.
\textbf{18.} $d^9$.
\textbf{19.} El acuocomplejo absorbe en el rojo y el naranja por una \omterm{def:b2:ligand-field:d-d}{transición d--d}: la
luz transmitida es azul.
\textbf{20.} Se vuelve blanco: sin ligandos agua alrededor del cobre (el sulfato ocupa su
lugar), el campo es más débil y la banda pasa al infrarrojo.
\textbf{21.} El agua vuelve a unirse al cobre y se forma de nuevo el acuocomplejo azul.
\textbf{22.} A una longitud de onda más corta: el amoníaco, por encima del agua en la serie, da un
desdoblamiento mayor (el azul violáceo intenso del complejo amminado).
\textbf{23.} No: $t_{2g}^6e_g^3$ es la única configuración.
\textbf{24.} Sí: un electrón desapareado.
\textbf{25.} $\boldsymbol{\Delta_o \approx \qty{1.8e4}{cm^{-1}}}$, unos
$\boldsymbol{\qty{215}{kJ/mol}}$.
\end{solution}
